\section{Introduction}\label{sec:intro}

% ═══════════════════════════════════════════════════════════════════
% 1.1  Breath as an affective sensory modality  (~450 words)
% Job: Establish breathing as special — dual nature, affective
%      richness, bidirectionality, therapeutic breathwork.
%      ADD: Slow breathing in contemplative/clinical contexts,
%      resonance breathing, breathwork and ASCs.
%      End with bridge sentence toward VR.
% ═══════════════════════════════════════════════════════════════════

\subsection{Breath as an Affective Sensory Modality}
Breathing carries both interoceptive and exteroceptive sensory information. Its
interoceptive component provides information about the body's internal
respiratory state, including lung inflation, airway and respiratory-muscle
activity, and oxygen and carbon dioxide levels
\citep{chan2024respiratory, engelen2023interoceptive,
weng2021interventions}. Because inhaled air stems from the external
environment, breathing also samples exteroceptive information about ambient
conditions, including odours, airflow, and air temperature, through its
coupling with orthonasal olfaction and thermo- and mechanosensory receptors in
the upper airways \citep{mori2022stress, mori2022processing,
zhou2011sensitivity, sabnis2008human, kytikova2021thermosensory}. Breathing is
also closely tied to affective experience. People sigh with relief after stress
\citep{Balban2023}, hold their breath during freeze responses
\citep{boiten1994emotions}, and hyperventilate during fear or panic
\citep{jerath2020respiratory, kreibig2010autonomic}. Arousing negative emotions
such as fear and anger are linked to shallower, more rapid breathing, whereas
relaxation and contentment correspond to slower, deeper respiratory patterns
\citep{boiten1994emotions, homma2008breathing, jerath2020respiratory}. Unlike
most visceral signals such as heartbeat, breathing is also amenable to
voluntary control \citep{sarkar2017functional}, making this relationship
bidirectional. Voluntarily changing respiratory patterns can modulate affective
state \citep{jerath2020respiratory, ashhad2022breathing}, a principle that
underlies breathwork as a therapeutic modality \citep{fincham2023effect,
Balban2023, Siebieszuk2025}. Breathwork techniques range from slow-paced
protocols that target relaxation, stress regulation, and improvements in heart
rate variability through sustained cardiovascular-respiratory coupling
\citep{Balban2023, Ahmed2021, Siebieszuk2025}, to high-ventilation protocols
such as cyclic hyperventilation that can facilitate more pronounced alterations
in affect and consciousness \citep{bahi2024effects, lewis2024breathwork,
havenith2025decreased}, effects increasingly posited as therapeutic, for
instance in facilitating emotional breakthroughs \citep{fincham2026airways,
fincham2024effects}. Breathwork is already established as a therapeutic
modality for stress, anxiety, and depression \citep{fincham2023effect,
hopper2023breathing, tymko2023breathwork}, and a growing number of virtual-reality applications also incorporate breathing for therapeutic and health-related purposes \citep{CortezVazquez2024, Pancini2025}. Yet this literature has so far explored only a relatively limited portion of the broader design space that breathing affords as an interaction modality in VR.

% ═══════════════════════════════════════════════════════════════════
% 1.2  Breath-based interactions in VR  (~600 words)
% Job: Survey the field with VR-specific texture, diagnose mixed
%      results, introduce design patterns and interaction modes,
%      identify gaps in both design and outcome space.
% ═══════════════════════════════════════════════════════════════════
\subsection{Breath-Based Interactions in VR}
Existing breath-based VR systems span meditation \citep{prpa2020articulating, Shamekhi2018}, anxiety regulation \citep{Bossenbroek2020}, stress management through heart-rate-variability biofeedback \citep{blum2019heart, rockstroh2019virtual}, and rehabilitation \citep{Shi2023, Alhammad2024}. In the context of our framework, we distinguish between two motifs in the design space regarding the mapping of the user's respiratory signal onto discrete visual objects, in contrast to environmental mappings. Object-oriented approaches typically place respiratory feedback in a discrete visual target. Third-person variants let the user watch an external entity respond to breathing, such as a cloud that approaches during inhalation and recedes during exhalation \citep{Tinga2019}. In first-person variants, the mapping is applied to a virtual body seen from the user's own vantage point. The ``embreathment'' illusion demonstrated that synchronising a virtual avatar's chest movements with the user's breathing enhances body ownership and agency \citep{MontiEtAl2020}, and mapping respiratory signals onto the expansion and contraction of a virtual avatar bubble viewed from a third-person perspective has been found to enhance agency and body awareness \citep{Wald2023}. What these object-oriented designs share is that respiratory feedback is concentrated in a single visual target, whether an external creature or a self-representation, which the user attends to as an indicator of their breathing. These designs are typically evaluated in terms of embodiment, measuring whether the temporal coupling between breathing and visual response enhances body ownership, felt agency, and awareness of the respiratory signal \citep{MontiEtAl2020, Wald2023}.

We contrast these with designs that embed respiratory feedback in the surrounding virtual environment rather than in a discrete object. DEEP \citep{Bossenbroek2020, vanRooij2016} and Flowborne \citep{rockstroh2021mobile} use diaphragmatic breathing to drive locomotion through immersive landscapes, rewarding paced breathing with movement and environmental change \citep{Bossenbroek2020, rockstroh2021mobile}. In a seated variant, \citet{blum2020development} mapped exhalation onto colour changes in flowers, trees, and rocks, keeping the feedback environmental but locally anchored to scene objects. A further group of environmental designs couples not the breathing signal directly but a slower downstream variable, typically heart rate variability, to ambient scene properties. In \citet{blum2019heart} and \textit{Heart Garden} \citep{Rook2026}, increasing coherence or heart rate variability gradually transforms the scene. These designs create calming feedback loops, but because the environment changes through a slower physiological proxy rather than respiration itself, moment-to-moment respiratory agency is attenuated. The locomotion-based designs provide immediate agency, but mainly as a progression mechanic rather than as a way of making surrounding space itself feel breath-coupled. In our design, we sought to synthesize these two approaches by mapping the user's breathing state onto an environment that expands and contracts with respiration. This provides immediate feedback and a sense of agency while presenting breathing as an ambient property of the surrounding space rather than as a distinct object requiring continuous attention. Specifically, the environment is composed of particles that respond to the user's breath such that inhalation expands the space outward and exhalation contracts it inward. In this way, changes in internal lung volume are mapped onto changes in external space. This mapping is closely related to what \citet{GoralEtAl2024} describe as exteroceptive--interoceptive sensory substitution, in which internal bodily signals are translated into external sensory events that remain experientially linked to the body rather than appearing as separate readouts. By mirroring the felt expansion and contraction of the thorax, our system was intended to preserve the natural structure of breathing experience and thereby function as an intuitive interface that can be understood with minimal explicit instruction \citep{Antle2009}. This choice was also informed by evidence from embodied cognition showing that people judge quantities as larger after inhalation than after exhalation \citep{belli2021breathink}, which may indicate a possible, though still speculative, link between respiratory state and the perception of space. More broadly, this embodied metaphor appears across a range of immersive biofeedback systems that map respiration onto expanding and contracting visual stimuli \citep{Wald2023, GoralEtAl2024, barton2024restorative, StepanovaEtAl2020}. Two existing systems are conceptually closest to the approach presented here. \citet{barton2024restorative} immersed participants in an abstract fractal-like environment and used a central breathing guide whose size participants controlled by spreading and closing their arms in time with the guide; the interaction was therefore structured as an explicit movement-synchronisation task rather than as an environment that responded directly to respiration. \citet{GoralEtAl2024} introduced exteroceptive--interoceptive sensory substitution in a system that translated respiratory signals into dynamic visual and auditory stimuli intended to feel continuous with the breathing body; however, their implementation externalized the signal through projection rather than embedding it in an immersive three-dimensional environment. Our system was designed to combine the strengths of these two approaches by pairing real-time environmental responsiveness that affords immediate respiratory agency with feedback distributed across surrounding space rather than concentrated in a single object or task. Rather than relying on semantic cues such as naturalistic scenery, skyboxes, or rendered objects, we pursued a bottom-up design in which the sense of surrounding space emerges from relations among many simple particles, drawing on Gestalt principles such as proximity, common fate, and continuity \citep{Wagemans2012} so that their dynamic arrangement, and specifically their responses to expansion and contraction, conveys depth and closeness dynamically within an abstract particle-based environment. This design rationale also motivates the broader question of whether coupling an internal bodily signal such as respiration to the surrounding sense of space can also reshape the sense of self, which we consider next through a review of relevant VR experiments.

\subsection{VR interventions testing the malleability of the self}
Virtual reality has been used extensively to investigate the malleability of bodily self-experience through multisensory manipulation. In full-body illusion paradigms, synchronous visuotactile stimulation between a participant's physical body and a seen virtual body can shift perceived self-location and induce a sense of ownership over the virtual body \citep{Lenggenhager2007, Petkova2008}. Later work showed that interoceptive signals can serve a similar synchronising function. Presenting a virtual body or silhouette that flashes in synchrony with the participant's heartbeat enhances body ownership and alters self-location even in the absence of tactile stimulation \citep{Aspell2013, Suzuki2013}. Respiratory signals have been used in a similar way. For example, synchronising a virtual avatar's chest movements, viewed from a first-person perspective, with the user's breathing enhances body ownership \citep{MontiEtAl2020}. Similarly, mapping respiratory signals onto the expansion and contraction of a virtual bubble, viewed from a third-person perspective, enhances agency and body awareness \citep{Wald2023}. Collectively, these studies show that a range of sensory modalities, including touch, breathing, and cardiac activity, can be used to reinforce visual feedback and thereby strengthen the sense of embodiment toward a virtual object or avatar from both first-person and third-person perspectives \citep{Lenggenhager2007, Petkova2008, Aspell2013, Suzuki2013, MontiEtAl2020, Wald2023}.

A related but distinct line of research concerns peripersonal space, the region immediately surrounding the body where the brain integrates multisensory signals in body-centered coordinates to maintain a dynamic boundary between self and environment \citep{Serino2019, Bufacchi2018}. This boundary is not fixed; it extends with tool use \citep{Maravita2003}, shifts with emotional valence and social context \citep{Bufacchi2018, Serino2019}, and interacts with interoceptive aspects of bodily self-representation \citep{Tsakiris2011, Azzalini2019}. Recent virtual-reality work further suggests that embodiment and multisensory manipulations can reshape peripersonal space in VR \citep{Petrizzo2024, Karakoc2025}. While the embodiment paradigms reviewed above manipulate the felt relationship between self and a visual body representation, peripersonal space research highlights that self-experience also depends on how the brain represents the spatial relationship between the body and the surrounding environment \citep{Serino2019, Bufacchi2018}.

Notably absent from both lines of research are paradigms that intervene through the surrounding environment rather than through a representation of the body itself. In most VR studies of embodiment and self-location, self-related experience is altered by manipulating how the body is represented, with any effects on perceived space arising indirectly through changes in how the body is seen in space \citep{Guy2023}. Although VR research has also developed techniques that manipulate space itself, such as scaling, translation, and redirected movement, these are typically used to support navigation and interaction rather than to investigate self-experience \citep{abtahi2022beyond}. By contrast, the present approach couples respiration not to a virtual body or discrete object, but to the spatial extent of the surrounding environment itself. This may engage the boundary between body-centred and environment-centred space in a way that body-focused paradigms do not.

There is a broad consensus that the sense of self is multidimensional \citep{milliere2020varieties, taves2020mystical}, and the present study assessed three complementary pictographic measures. These were body-boundary salience \citep{Dambrun2016}, spatial frame of reference \citep{Hanley2019}, and perceived self-size \citep{vidal2025inducing}. At the same time, the intervention was expected to affect not only self-experience, but potentially the broader structure of subjective experience. The present study therefore also considered its effects in relation to psychometric frameworks from altered-states research, which informed both measurement and the design rationale of the VR system described next.


% ═══════════════════════════════════════════════════════════════════
% 1.4  Breathwork and altered states of consciousness  (~400 words)
% ═══════════════════════════════════════════════════════════════════
\subsection{Breathwork as an experiential design space for virtual reality}

Embedding a breathing protocol in a responsive virtual environment introduces specific design constraints. Different breathing practices produce different respiratory patterns, including variations in the pace, duration, and force of inhalation and exhalation. Visual feedback coupled to these dynamics must therefore avoid becoming dizzying, claustrophobic, disorienting, or overwhelming. In VR, large-amplitude radial scene motion, particularly in the visual periphery, is a documented source of cybersickness \citep{Chang2020vrsickness, Poehlmann2021vection}, and contracting environments can elicit involuntary defensive postural responses even in non-phobic users \citep{KodithuwakkuArachchige2022walls}. These risks are compounded when the user is in a heightened or unfamiliar experiential state, given that negative affect and arousal are associated with increased VR sickness severity \citep{Kaufeld2022emotions}. The choice of breathing protocol therefore shapes not only what physiological and experiential effects to expect, but also what the visual environment can safely and comfortably do in response.

High-ventilation breathwork practices such as holotropic breathing, conscious connected breathing, and cyclic hyperventilation involve sustained increases in breathing rate and depth that, over extended sessions, induce pronounced physiological changes, including respiratory alkalosis and reductions in cerebral blood flow, and can give rise to marked changes in subjective experience \citep{lewis2024breathwork, Kartar2025, fincham2024effects}. Participants report experiences of oceanic boundlessness, in which the felt boundary between self and world dissolves into a sense of unity or merging with the surroundings; visionary restructuralisation, involving vivid geometric imagery, complex visual scenes, and synesthetic blending of sensory modalities; and dread of ego dissolution, characterised by a felt loss of control over one's thoughts and sense of identity \citep{bahi2024effects, havenith2025decreased, lewis2024breathwork}. These experiential dimensions overlap with those reported under medium- to high-dose psilocybin as measured by the 11D-ASC \citep{studerus2010, havenith2025decreased}, and are often accompanied by emotional catharsis and strong bodily sensations such as tingling, dizziness, and autonomic arousal \citep{havenith2025decreased, lewis2024breathwork, fincham2024effects, fincham2026airways}. However, the rapid, large-amplitude chest excursions characteristic of high-ventilation protocols would produce correspondingly rapid expansion and contraction of a breath-coupled particle environment, and the accompanying physiological arousal and altered experiential state would coincide with precisely the visual conditions most likely to provoke cybersickness and spatial discomfort. For these reasons, we chose to facilitate slow-paced breathing rather than high-ventilation breathwork in the present system.

Slow-paced breathing occupies a different region of this subjective and affective space. Its established effect profile is characterised by relaxation, reduced anxiety, increased comfort, and calm alertness, supported by parasympathetic dominance and enhanced heart rate variability \citep{zaccaro2018breath, Balban2023, fincham2023effect}. Slow breathing is also one of the most common attentional anchors in guided meditation practice \citep{Gerritsen2018, zaccaro2018breath}, and evidence from experienced pranayama practitioners indicates that very slow nasal breathing can induce altered states of consciousness characterised by altered body perception, distortions in time sense, unusual salience, increased joy, and reduced anxiety \citep{zaccaro2022neural}. The mild but documented non-ordinary experiential profile of slow-paced breathing provided a low baseline against which any augmentation by the visual environment would be easier to detect, while keeping the environment's spatial dynamics within a comfortable perceptual range.

Augmenting subjective experience is clinically relevant because experiential intensity during altered-state interventions predicts therapeutic outcomes \citep{Roseman2018, DosSantos2025, havenith2025decreased}. We chose a slow-paced protocol precisely because its mild baseline provides sensitivity for detecting whether we can enhance specific attributes of the subjective experience through VR. To characterise this experiential profile and compare it with existing breathwork and psychedelic literatures, we adopted the 11D-ASC \citep{studerus2010, Dittrich1998}, the standard instrument across these fields. Its subscales capture the phenomenological features (e.g., boundary dissolution, changed imagery, synesthetic blending) that a breath-coupled environment might selectively amplify, and can be summarised into three higher-order dimensions reflecting positive, distressing, and perceptual effects \citep{Stocker2025}. Because we also examine how these dimensions evolve over the course of a breathing session, we include psychometric time-series ratings, which we refer to as temporal experience tracers \citep{Jachs2022}. We match the tracers one-to-one with the 11D-ASC dimensions to compare global ratings with time-resolved intensity profiles of the same experiential constructs. The motivation to enhance subjective experience is also reflected in the aesthetic design of the visual environment, which we describe next.


 
% ═══════════════════════════════════════════════════════════════════
% 1.5  Particle dynamics as an affective bioresponsive design medium
% ═══════════════════════════════════════════════════════════════════

%P0
\subsection{Particle dynamics as a bioresponsive design medium}
Viscereality is a bioresponsive VR system developed through an iterative design process, parts of which have been reported previously with respect to its design rationale, technical implementation, and interaction mechanics in human-computer interaction contexts \citep{fejer2025viscereality, fejer2026breathing}. The present study provides its first controlled empirical evaluation, testing whether the system's particle dynamics measurably shape subjective experience during breathwork. As described above, the system maps lung volume onto a surrounding particle field through interoceptive-exteroceptive sensory substitution \citep{GoralEtAl2024}. Inhalation expands the environment outward, whereas exhalation contracts it inward. Existing implementations of this type of coupling typically use modelled surfaces, either presented as localized targets linked to the user \citep{Wald2023} or projected onto the walls of a physical room \citep{GoralEtAl2024}. In both cases, the user watches a representation change size, but a uniform surface that grows or shrinks does not, on its own, strongly convey that the surrounding space itself is approaching or receding; that impression depends on motion-based depth cues that arise when many elements move relative to one another \citep{Rogers1979}. We designed Viscereality as a particle system specifically to afford this. The environment is composed of approximately 2,500 particles distributed on a geodesic icosphere, and the sense of spatial expansion and contraction arises from Gestalt grouping among these elements, specifically proximity, common fate, and continuity \citep{Wagemans2012}, rather than from a texture mapped onto a surface. Because each particle independently carries multiple visual properties (position, size, colour, transparency, motion trajectory), the same elements that produce the breath-coupled spatial dynamics can simultaneously encode additional biofeedback channels, such as heart rate variability coherence \citep{fejer2026breathing}, without requiring separate visual layers. And because inter-particle interaction rules, specifically oscillator coupling across neighbours, generate emergent aesthetic properties such as dynamic symmetry and geometric coherence, these can be parameterised independently of the breathing mechanics, making the aesthetic texture of the environment an experimentally controllable variable. Throughout the breathing cycle, the particle field remains continuously active. Surface-wave oscillations, chromatic cycles, and shape fluctuations continue during inhalations, exhalations, and both breath-hold phases. The within-VR manipulation therefore does not switch visual motion on or off; rather, during breath holds the high-symmetry condition transiently aligns these ongoing particle cycles into coherent geometric organisations, whereas the asymmetric condition leaves the same dynamics desynchronised. To our knowledge, no existing interoceptive-exteroceptive sensory substitution system takes this particle-based approach. Aesthetics in VR are often treated as incidental to intervention design, yet recent work identifies them as among the strongest predictors of emotional response and experience quality in immersive environments \citep{Marcolin2021, Cheiran2025}. In Viscereality, aesthetics are not a cosmetic layer but emerge from the same particle dynamics that carry the respiratory coupling, drawing on a longer history of expressive and aesthetic design in particle-based systems applied in visual effects animation.

% ── P1: Particle systems as a generative medium ─────────────────
\subsubsection{Particle systems as a generative medium}
Particle systems were introduced by \citet{Reeves1983} to represent visual phenomena that have no fixed boundary and no stable shape, including fire, smoke, clouds, and water. In his formulation, each particle is born with stochastically assigned properties (position, velocity, colour, size, transparency, lifetime), moves and changes over time, and eventually vanishes. The overall form is never fully specified; it emerges from the aggregate of many simple elements, each governed by a small set of controllable parameters. This is the property that makes particle systems useful as a bioresponsive medium because each particle independently carries multiple visual channels, so a single population of elements can simultaneously respond to physiological signals through spatial position, colour shifts, size, transparency, or motion trajectory without requiring separate visual layers for each mapping. Building on this work, \citet{Reynolds1987} introduced interaction between elements, enabling coordinated collective behaviour to emerge from local neighbour rules, whereas in a standard particle system each element responds only to global parameters and physics without awareness of its neighbours. Reynolds specified three local rules. They were to avoid collision with nearby elements, match their velocity, and maintain proximity to them. This became foundational for the ability of 3D animation to produce the fluid, coordinated group motion recognisable in flocking birds, schooling fish, or stampeding herds, all arising bottom-up from neighbour interactions without a centrally scripted choreography.


% ── P2: Particle systems as an affective medium ─────────────────
\subsubsection{Particle systems as an affective medium}
Research on abstract emotion representation shows that affective meaning can be conveyed through low-level perceptual parameters of abstract expressions \citep{deRooij2013}. In particular, affective valence can be mapped through semantically minimal features such as motion speed and trajectory smoothness \citep{Bartram2010, Pollick2001, Burger2020}, angular or rounded shape contours \citep{Blazhenkova2018, Aronoff2006}, and line orientation \citep{Damiano2023}. In conceptual terms, \citet{Khatin-Zadeh2019} argue that motion provides an effective vehicle for representing emotion because both unfold dynamically over time, allowing changes in affective state to be expressed through movement, direction, and position in space. This claim has also been borne out empirically in swarm robotics, where the collective motion of distributed abstract disc-shaped agents has been shown to alter the observer’s state of flow and sense of time relative to observing a single moving agent \citep{Kaduk2024}. \citet{Dietz2017} found that higher speed in swarm motion increased perceived emotional intensity, while smooth trajectories, relative to jittery motion, were associated with more positive affect. Extending this approach, \citet{Santos2021} translated trajectory and motion into swarm choreographies for Ekman’s six basic emotions \citep{ekman1993facial}. Rounded and smooth configurations marked positive or low-arousal states, while angular, clustered, or dispersive formations with fast or slow trajectories marked negative and higher-arousal states. Participants identified the intended emotions above chance \citep{Santos2021}. What makes these findings from swarm robotics relevant here is that their expressive effects arise from relational motion among many simple elements rather than from the form of any individual agent. The same principle holds in \textit{Viscereality}, where local coupling among particles governs trajectory speed, smoothness, synchrony, clustering, and dispersal, producing higher-order collective motion patterns from which affective qualities may emerge. The particle field can therefore be treated as a medium for conveying affective states through abstract visual dynamics rather than semantic or figurative cues.


% ── P3: Particle Systems as an aesthetic medium ───────────────────────────

\subsubsection{Particle systems as an aesthetic medium}

Beyond serving as a medium for affective mapping, we designed the particle field as an aesthetic means of rendering affective states through abstract spatiotemporal organization. Research on abstract visual art suggests that non-figurative forms evoke aesthetic and affective responses through the isolation of simple features, the grouping of many elements into emergent wholes, and the amplification of perceptual qualities through Gestalt grouping principles \citep{deRooij2013, Ramachandran1999}. In our particle system, these same grouping principles also help convey depth when approximately 2,500 ring-shaped particles expand and contract around the user. Because each particle is visually minimal, higher-order forms are not pre-given but emerge through local relations among neighbouring elements, allowing circles, lattices, polygons, and other geometric organisations to arise through changes in three-dimensional location. The appearance of such forms can itself become aesthetically salient, as the resolution of an initially indeterminate field into a coherent Gestalt is known to produce pleasurable moments of insight or `aesthetic Aha' moments \citep{muth2013aesthetic, spehar2017koffka, spee2024dynamics}. In Viscereality, this made it possible to translate oscillator-level control into changes in overall visual order, symmetry, and coherence without moving to figurative imagery.

Another design influence was the geometric and kaleidoscopic qualities of hallucinatory and psychedelic imagery, which are often marked by geometric patterns, bilateral symmetry, and perceptual ambiguity \citep{Kluver1966, Billock2012, Lawrence2022, Rubin2010}. This motivated our use of ring-shaped particles, as their overlaps can give rise to distinctive emergent patterns during the breathing cycle. Because these qualities emerge from shifting relations among many simple particles rather than from fixed symbolic forms, they remain open in meaning and foreground a defining affordance of the medium, namely its capacity to continually reorganise and transform its aesthetic structure over time.

In the context of VR design, \citet{Glowacki2024} terms such non-figurative scenes a weak representational aesthetic, in contrast to photorealistic depictions that render scenes and bodies as recognizable real-world entities. In practice, this aesthetic has been applied in VR scenarios designed to induce altered states of consciousness through the representation of participants' bodies as luminous, semi-corporeal avatar bodies composed of particle clouds with mergeable body boundaries, contributing to the malleability of self-experience and altered subjective states \citep{Glowacki2020, Glowacki2022, Glowacki2024, vidal2025inducing}. Our approach adopts a related design aesthetic while focusing more specifically on diminishing the felt boundary between self and environment.

To evaluate the particle system's capacity to manipulate aesthetic properties, we operationalized symmetry as an additional independent variable that varied across conditions. Previous studies have linked visual symmetry to more positive affect and perceptual fluency \citep{Bertamini2013, Pecchinenda2014, Makin2012}, while synchrony in group movement has been linked to greater enjoyment and more positive affective appraisal \citep{Vicary2017, Mogan2017}. We therefore asked whether higher degrees of symmetry would enhance positively valenced experiences in the context of Viscereality \citep{johnson2023qualia}. We were also interested in whether participants would be able to detect and discriminate between the two conditions.
 

\subsection{Present Study}

The present study provides the first controlled empirical evaluation of Viscereality with naive participants. We used a preregistered within-subjects design with $N = 39$. Each participant completed three 10-minute audio-guided box-breathing sessions in counterbalanced order. These were a high-symmetry condition, in which positive oscillator coupling transiently aligned ongoing particle oscillations into dynamic geometric coherence during breath-hold phases; an anti-symmetry condition, in which negative coupling produced a visually irregular particle field throughout; and a dark-screen control (see Section~\ref{subsubsec:breathing_indicator}). The two particle conditions used identical particles, breath-coupled spatial dynamics, colour palette, and audio guidance, differing only in oscillator coupling parameters. The central within-VR contrast therefore concerned the aesthetic organisation of an otherwise identical breath-responsive scene.

We assessed four substantive outcome domains together with a methodological blinding assessment. Breathing adherence was quantified as the phase-locking value between the observed respiratory signal and an ideal reference trajectory. Self-related experience was assessed with three pictographic measures, namely body-boundary salience \citep{Dambrun2016}, spatial frame of reference \citep{Hanley2019}, and perceived self-size \citep{vidal2025inducing}. Altered states of consciousness were assessed retrospectively with the 11D-ASC \citep{studerus2010, Dittrich1998}. Blinding efficacy between the two active VR conditions was assessed retrospectively with a two-stage procedure. Before unmasking, participants completed 12 probe questions covering condition-specific, shared, and sham features. After being told that the two active scenes differed, they indicated in which order the high-symmetry and anti-symmetry scenes had been presented, together with confidence ratings. These responses were used to estimate condition-identification blinding with a standard single-group Bang blinding index and an exploratory confidence-weighted directional variant \citep{Bang2004}.

We hypothesised that, relative to the dark-screen control, the bioresponsive VR conditions would increase breathing adherence and altered-state intensity and would also influence the three preregistered self-related measures. Within VR, we further expected the high-symmetry condition to enhance positively valenced experiences and produce stronger alterations in self-related experience than the anti-symmetry condition. The preregistration did not specify a separate directional prediction for each pictographic outcome. We therefore treated the three outcomes as one family and controlled the false-discovery rate across them.


